En el articulo los autores desarrollaron un modelo probabilístico
para el pronóstico de demanda eléctrica del sistema de Baja
California Sur basado en la combinación de redes neuronales
convolucionales y regresión cuantílica (CNN-QR), complementado
mediante la Rainbow Technique for Categorical Features (RTCF).

La regresión cuantílica permite estimar diferentes regiones de la
distribución condicional de la demanda mediante

\begin{equation}
Q_{Y|X}(\tau|X)
=
X^{T}\beta_{\tau},
\end{equation}

donde $\tau \in (0,1)$ representa el cuantil considerado.

El entrenamiento puede expresarse mediante la función de pérdida
cuantílica

\begin{equation}
\rho_{\tau}(u)
=
u
\left[
\tau-\mathbb{I}(u<0)
\right],
\end{equation}

donde

\begin{equation}
u=y-\hat{y}.
\end{equation}

De manera equivalente,

\begin{equation}
L_{\tau}(y,\hat{y})
=
\begin{cases}
\tau(y-\hat{y}), & y\geq\hat{y},\\
(1-\tau)(\hat{y}-y), & y<\hat{y}.
\end{cases}
\end{equation}


La arquitectura propuesta incorpora una capa convolucional
unidimensional con 128 filtros y un kernel de tamaño 3, seguida de
MaxPooling, capas densas, regularización $L_2$ y Dropout. Además,
RTCF combina diferentes variables categóricas y temporales para
construir una nueva característica denominada \textit{Category}.

Entre las principales ventajas del método se encuentra su capacidad
para producir pronósticos probabilísticos y representar relaciones
no lineales dentro de los datos. Asimismo, la integración de RTCF
permite incorporar de manera conjunta información relacionada con
el día, mes, estación, días festivos y hora del día.

Los resultados muestran que CNN-QR obtiene menores pérdidas
cuantílicas que la regresión cuantílica convencional. Por ejemplo,
para los cuantiles $0.1$, $0.5$ y $0.9$, CNN-QR obtuvo pérdidas de
$8.84$, $11.08$ y $8.70$, respectivamente, mientras que la regresión
cuantílica obtuvo valores de $39.00$, $21.50$ y $34.34$.

No obstante, el método presenta una mayor complejidad computacional
que los modelos estadísticos tradicionales debido al entrenamiento
de una red neuronal profunda. También presenta una menor
interpretabilidad debido a las transformaciones internas realizadas
por las capas convolucionales y densas.

Por otra parte, RTCF produce una cantidad considerable de
combinaciones categóricas, algunas de las cuales aparecen con muy
poca frecuencia en el conjunto de entrenamiento. Esto puede
dificultar la generalización del modelo hacia combinaciones que no
estén suficientemente representadas.

La principal limitación del estudio es su alcance geográfico, dado
que el modelo fue desarrollado y evaluado exclusivamente con
información correspondiente a Baja California Sur. El conjunto de
datos comprende aproximadamente 15\,336 observaciones horarias
registradas principalmente entre 2019 y 2020, por lo que no se
demuestra su capacidad de generalización a otras regiones del
Sistema Eléctrico Nacional.